\documentclass{article}

\usepackage{amsmath}
\usepackage{url}
\usepackage{booktabs}
\usepackage{epsfig}
\usepackage{graphics}
\usepackage{latexsym}
\usepackage{amssymb}
\usepackage{comment}
\usepackage{fancyvrb}
\usepackage{moreverb}
\usepackage{mathptmx}
\usepackage[T1]{fontenc}
\usepackage[utf8]{inputenc}
\usepackage{latexsym,color}
\usepackage[french]{babel}
\usepackage{pstricks,pst-node,pst-text,pst-3d,pst-tree,pst-grad}
\usepackage{xcolor}


\newcommand{\fig}[2]
{
 \begin{figure*}[hbt]
  \begin{center}
  \includegraphics[width=10cm]{#1}
  \caption{\label{#1}  #2}
  \end{center}
 \end{figure*}
}

% small et verbatim: to be coherent with \file
\newenvironment{smallverbatim}%
{\endgraf\small\verbatimtab}{\endverbatimtab}


\definecolor{mygray}{gray}{0.90}
\definecolor{mygreen}{rgb}{0,0.5,0}
\definecolor{maroon}{rgb}{0.5,0,0}
\usepackage{listings}


\title{Questions et exercices sur l'indexation par partitionnement}
\author{Philippe Rigaux}

\date{\today}

\begin{document}

\maketitle

\textbf{Indiquez votre nom}\,:


\begin{enumerate}

 \section{Distribution par partitionnement}

  \item \textbf{Expliquez brièvement le rôle du routeur}
  
  \vspace*{4cm}
  
  \item \textbf{Quelles opérations ne sollicitent qu'un seul serveur\,?}
  
  \vspace*{4cm}
  
\section{Partitionnement par intervalle}
  
  
 \item \textbf{Peut-on créer un index non-dense sur une collection non-triée\,?}
 
   \vspace*{4cm}
  
  
 \item \textbf{Pourquoi indexer les fragments, et pas les documents eux-mêmes\,?}
 
   \vspace*{4cm}
 
 \item \textbf{La clé de ma collection de vidéos est le titre. Est-il possible de faire
      une recherche par année dans un système distribué\,?}
 
   \vspace*{4cm}
    
 \item \textbf{Un même serveur peut-il gérer plusieurs fragments\,?}
   
   \vspace*{4cm}
      
 \item \textbf{Est-il nécessaire de trier un fragment\,?}
   
   \vspace*{4cm}
   
 \item \textbf{Décrivez brièvement l'opération de split.}
   
   \vspace*{4cm}
    
    \section{Hachage cohérent}
    
    \item Si on prend une fonction de hachage basée sur la fonction modulo $N$, où $N$ est le nombre 
       de serveurs, que se passe-t-il si on ajoute des serveurs sans changer la fonction\,?
       
          \vspace*{4cm}
       
       \item Si on change la fonction en (modulo $N+1$), que se passe-t-il pour les données déjà insérées\,?
       
          \vspace*{4cm}
       
       \item Dans un structure Chord, exprimez la distance entre  \emph{friends}$[i]$ et
  \emph{friends}$[i+1]$ sur l'anneau de hachage, et montrez qu'une opération de recherche coinverge
      en temps logarithmique.
      
         \vspace*{4cm}

\begin{figure*}[hbt]
  \begin{center}
   \includegraphics[width=8cm]{images/chord-friends}
  \caption{\label{figex:chord}  Un anneau Chord}
  \end{center}
 \end{figure*}
       
      \item Soit l'anneau Chord de la figure~\ref{figex:chord}.Quels sont les amis
         du pair 
$p11$, situé à la position 3? 
        \vspace*{4cm}

\item Développez un exemple d'une recherche dans Chord avec $m=10$, dans le pire des cas.

\end{enumerate}


\end{document}